% ─────────────────────────────────────────────────────────────────────
% Раздел 1 — Задача
% ─────────────────────────────────────────────────────────────────────

\section{Задача}

Представьте себе подземную угольную шахту: километры тоннелей на
нескольких горизонтах, в среднем по сорок-сто шахтёров на смену, у
каждого фонарь, который раз в пятнадцать секунд шлёт радиопакет с
показаниями нескольких ближайших маяков связи, а раз в десять секунд
--- отдельный пакет с показаниями газовых сенсоров. Эти пакеты летят
в диспетчерскую через смесь UDP-радио и LTE-шлюзов, доходят с
задержками, иногда дублируются, иногда теряются.

В диспетчерской на стене висит экран с картой шахты и состоянием
смены: где находится каждый шахтёр, у кого тревога по газу, кто
прижал кнопку SOS, у кого вот-вот сядет батарейка фонаря. Этот экран
обслуживает наш сервис.

\subsection{Что должен делать сервис}

Если разложить требования \emph{диспетчера} на конкретные пайплайны,
получается три независимых подсистемы:

\begin{description}
\item[Контроль состояния шахтёра.] Из RSSI-пакета извлекаем не
показания маяков (это работа другой подсистемы), а \emph{косвенные
факты}: пакет вообще пришёл, в нём были показания, фонарь двигался,
напряжение батареи в норме. По каждому из этих фактов держим
независимый таймер и эмитим соответствующий алерт когда таймер
перевалил порог: ``не было пакетов больше двух минут'',
``пакеты идут, но маяков не слышит больше пяти'', ``не двигался дольше
порога'', ``батарея ниже 3{,}5 В с дребезгом по гистерезису''.
Плюс импульс SOS (нажатая кнопка) --- мгновенный без таймера.

\item[Локация по маякам.] Для каждого шахтёра в скользящем
60-секундном окне с шагом 5 секунд берём топ-2 маяков с самой сильной
медианной RSSI и из них приближённо вычисляем координаты шахтёра
через интерполяцию между двумя маяками пропорционально расстоянию.
Окно скользит каждые 5 секунд, поэтому диспетчерский экран обновляется
не реже чем раз в 5 с; при опаздавшем пакете окно
\emph{переоткрывается} и эмитит \texttt{Retract} старого значения с
последующим обновлённым \texttt{Data} --- так картинка остаётся
консистентной.

\item[Газовые алерты с обогащением координатами.] Газовый пакет
приходит \emph{отдельно} и сам по себе не несёт ID маяков или
позиции. Когда газ превышает порог --- алерт идёт немедленно
(безопасность важнее красоты), но без координат, если их ещё не было.
Как только локация для этого шахтёра становится известна (хотя бы по
одному маяку), алерт автоматически \emph{обновляется}: старая версия
без координат retract'ится, новая с координатами эмитится. То же самое
при смене уровня опасности (Warning $\to$ Critical) или
существенном изменении ppm. Когда газ возвращается в норму --- алерт
снимается через \texttt{Gas\_resolved}.
\end{description}

Эти три подсистемы независимы по бизнес-логике, но обмениваются
данными: газовый пайплайн \emph{потребляет} результат пайплайна
локации, чтобы обогатить свои алерты координатами.

\subsection{Что в этой задаче сложного}

Если бы это была учебная задача из курса по базам данных, её можно
было бы решить за полчаса: батчевый ETL раз в минуту, JOIN по
шахтёру, отрисовка в дашборд. Но в реальной шахте такое решение
\emph{бесполезно}, и причин этому несколько.

\paragraph{Время событий не равно времени поступления.} Пакет от
шахтёра M1 с временной меткой \texttt{t=88s} может прийти в
диспетчерскую в момент \texttt{t=210s}: UDP в шахте теряет пакеты,
LTE-шлюз делает повторные попытки, на агрегаторе очереди.
Если мы будем считать алерты по wall-clock времени поступления, то
получим ``M1 не двигался последние 2 минуты'' в момент когда M1 уже
давно дома --- его пакет о движении просто опаздывает по дороге.
Сервис должен работать по \emph{event-time} (времени самого
события), а не по wall-clock.

\paragraph{Идемпотентность.} Тот же UDP-механизм даёт дубликаты:
пакет может прийти дважды, потому что шлюз ретранслировал его при
сомнениях. Если мы будем считать тревогу по второму дублю
``независимо'', получим \emph{два} алерта SOS на одно нажатие кнопки.
Диспетчер не должен видеть фантомных тревог.

\paragraph{Опоздавшие пакеты ломают окна.} Окно
\texttt{[60s, 120s)} закрылось в 120с, и пайплайн эмитнул финальное
top-2 маяков. Но в 145с пришёл опоздавший пакет с
\texttt{ts=90s} --- он \emph{принадлежит} уже закрытому окну. Нужно
либо его выбросить (теряем данные), либо переоткрыть окно, пересчитать
агрегат и эмитнуть новое значение, помечая старое как недействительное.
Второе правильно для диспетчера; первое было бы дешевле, но даёт
непредсказуемый дашборд (на одной и той же временной метке могут быть
разные значения в разные моменты wall-clock).

\paragraph{Газовый алерт без координат лучше чем нет алерта.}
Газ начинает выделяться \emph{раньше} чем шахтёр успевает прийти под
маяк --- особенно в забое, где маяки ставят редко. Если ждать
координат прежде чем эмитить тревогу о 120 ppm CO, можно
потерять минуты на ожидании. Алерт идёт сразу. Но как только
координаты появляются --- они должны \emph{подтянуться} к уже
существующей тревоге, не порождая дубль и не давая диспетчеру повод
закрыть старый алерт и открыть новый вручную.

\paragraph{Объём.} Шахта на сорок шахтёров --- мелочь. Реальные
большие шахты в Кузбассе или Караганде имеют до четырёх тысяч
работников в смену. RSSI-пакетов от такой шахты получается порядка
270 ev/s, газовых --- ещё 400 ev/s. Это
немного по меркам стримов, но \emph{пайплайн должен держать нагрузку
с запасом}: пиковые моменты (несколько горящих участков, грозовая
помеха в радио, перезапуск шлюза) дают взрывы в десятки тысяч событий
в секунду на короткий период. Сервис не должен в эти моменты
терять алерты.

\subsection{Что есть готового в miniFlink}

Библиотека даёт нам почти всё что нужно для этого сервиса, без
необходимости писать stream-runtime самостоятельно:

\begin{itemize}
\item \texttt{Pipe.event\_time \textasciitilde lateness} ---
вставка watermark'ов в поток по event-time;
\item \texttt{Pipe.dedup} --- дедупликация по rule-функции;
\item \texttt{Pipe.window\_agg\_keyed} с
\texttt{allowed\_lateness} --- скользящие окна с retract'ами при
опоздавших пакетах;
\item \texttt{Agg.median}, \texttt{Agg.group\_by},
\texttt{Agg.top\_k\_by} --- набор комбинируемых агрегатов;
\item \texttt{Pipe.process\_keyed} --- per-key state с self-timers
(для FSM состояния шахтёра);
\item \texttt{Mf\_event.union} --- объединение потоков с
выравниванием watermark'а по минимуму (для трёх-поточного газового
пайплайна);
\item \texttt{Pipe.materialize} --- сворачивание retract-богатого
потока в финальную таблицу записей (для sink-а).
\end{itemize}

Чего \emph{нет} --- это бизнес-логики и доменных типов. Их пишем мы;
структура ниже как раз и есть рассказ о том, как мы их написали.

\subsection{Что увидим дальше}

Статья идёт от внешнего наружу. В разделе~\ref{sec:inputs}
разбираем устройство входов: почему два независимых потока вместо
одного, какие в них аномалии. В разделе~\ref{sec:pipelines} --- три
пайплайна по очереди: FSM состояния, скользящие окна с медианой и
интерполяцией позиции, retract-обогащение газовых алертов.
Раздел~\ref{sec:retract} собирает наблюдения о retract-семантике как
о сквозном механизме, проходящем через все три пайплайна.
Раздел~\ref{sec:performance} приводит измерения производительности
(на машине автора-эксперимента и на пользовательском i7-8700 с OCaml
5) и описывает направленные оптимизации, которые мы делали.
Раздел~\ref{sec:conclusions} --- что отложено, что не делалось
намеренно, и при каких условиях стоит реализовать недостающее.
